% =====================================================================
%  Chapter 1 -- Introduction (limited to one page)
% =====================================================================
\chapter{Introduction}
\label{ch:Introduction}

The continued electrification of the global economy is driving a sustained
increase in electricity demand~\cite{international_energy_agency_electricity_2025}.
Power electronic converters (PECs) enable this shift by electrifying
transport~\cite{bosshard_multi-objective_2016} and energy
systems~\cite{steimer_enabled_2010}. The power density of PECs has risen for decades,
driven by cost, functionality and hard weight and space
limits~\cite{steimer_enabled_2010,biela_cooling_2007}, and is projected to keep
rising~\cite{morroni_understanding_2023}. As converters shrink, however, the
same losses must be removed from a smaller volume, so thermal management
increasingly limits how far power density can
rise~\cite{biela_cooling_2007,morroni_understanding_2023}.

Inductive power transfer (IPT) chargers for electric vehicles illustrate this
challenge. At the standard \SI{85}{\kilo\hertz}~\cite{society_of_automotive_engineers_r_2024},
a \SI{10}{\kilo\watt} system dissipates several hundred watts, much of it in pads
that generate strong alternating magnetic fields~\cite{kim_thermal_2022}. Many
commercial designs already actively cool the coils~\cite{lawton_power_2024}, and
Liu \emph{et al.}\ call for new technologies to cool wireless
chargers~\cite{liu_thermal_2025}. Water cooling, however, is limited by the low
thermal conductivity of water~\cite{laloya_heat_2016}. Galinstan, a liquid metal
with much higher thermal conductivity~\cite{noauthor_gainsn_nodate}, can instead
be driven by a magnetohydrodynamic (MHD) pump with no moving
parts~\cite{hodes_potential_2014}, but its low volumetric heat capacity and
electrical conductivity bring trade-offs~\cite{hodes_potential_2014,an_liquid_2024}.
The extent of its advantage therefore remains unclear, as few studies compare it
with water at converter scale~\cite{fan_ultra-efficient_2026}.

This project builds on the 2025 Part~IV Project~\#51, in which Kwon and Loh
built an MHD-pumped Galinstan cooling loop and recommended comparing it with
water cooling and testing it near magnetic
components~\cite{kwon_project_2025,loh_project_2025}. Its aim is to evaluate
MHD-pumped Galinstan as a coolant for high-power-density PECs, and to identify
where in an IPT system it is most useful. The objectives are to (i)~compare the
thermal resistance and pumping power of the Galinstan loop with those of a water
loop using the same cold plate~\cite{ekbote_project_2026}; (ii)~explain and
extend the measurements with thermal models and a conjugate heat transfer
simulation; and (iii)~use a power-dense \SI{85}{\kilo\hertz} inductor to
demonstrate Galinstan cooling of magnetics and measure the eddy-current loss
induced in the coolant. These objectives are framed as two research questions
in \cref{sec:research-questions}. \Cref{ch:literature} reviews the literature,
\cref{ch:experimental} describes the methodology, \cref{ch:results} presents
the results, \cref{sec:discussion} discusses them and \cref{ch:conclusion}
concludes the report.
